% Generated draft: requires mathematical and source review.
\documentclass[11pt]{article}
\usepackage[margin=25mm]{geometry}
\usepackage{fontspec}
\setmainfont{Latin Modern Roman}
\usepackage{amsmath,amssymb,mathtools,unicode-math}
\setmathfont{Latin Modern Math}
\usepackage{enumitem,needspace,xurl,hyperref}
\setlength{\parindent}{0pt}
\setlength{\parskip}{5pt}
\setlength{\emergencystretch}{3em}
\newcommand{\N}{\mathbb{N}}
\newcommand{\Z}{\mathbb{Z}}
\newcommand{\Q}{\mathbb{Q}}
\newcommand{\R}{\mathbb{R}}
\newcommand{\C}{\mathbb{C}}
\newcommand{\F}{\mathbb{F}}
\newcommand{\A}{\mathbb{A}}
\newcommand{\PP}{\mathbb{P}}
\newcommand{\E}{\mathbb{E}}
\newcommand{\eps}{\varepsilon}
\newcommand{\dd}{\,\mathrm{d}}
\begin{document}
\section*{UnsolvedMath ID: 30004181}
\textbf{Source problem code:} catalogue entry
\subsection{Definitions and mathematical statement}
Let \(\mathcal M\) denote the source's intended class of mildly dissipative disk diffeomorphisms, equipped with its intended topology \(\tau\). For \(f\in\mathcal M\), let \(h_{\mathrm{top}}(f)\) denote its topological entropy, and write \(\operatorname{IR}(f)\) for the property that \(f\) is infinitely renormalizable. Define
\[
\mathcal Z=\{f\in\mathcal M:h_{\mathrm{top}}(f)=0\},\qquad
\mathcal I=\{f\in\mathcal Z:\operatorname{IR}(f)\}.
\]
The boundary relative to \((\mathcal M,\tau)\) is
\[
\partial_{\mathcal M,\tau}\mathcal Z
=\overline{\mathcal Z}^{\,\mathcal M,\tau}
\setminus\operatorname{int}_{\mathcal M,\tau}(\mathcal Z).
\]
The source asks whether
\[
\partial_{\mathcal M,\tau}\mathcal Z=\mathcal I,
\]
or, equivalently, whether every \(f\in\mathcal M\) satisfies
\[
f\in\partial_{\mathcal M,\tau}\mathcal Z
\quad\Longleftrightarrow\quad
h_{\mathrm{top}}(f)=0\ \text{and}\ \operatorname{IR}(f).
\]
This is a conditional formulation: the supplied record does not specify the differentiability class, topology, disk and mapping boundary conventions, or definitions of mild dissipation and infinite renormalizability. These must be recovered from the cited sources before the displayed question determines a fully specified mathematical proposition. No assertion is made here for all dissipative disk diffeomorphisms.
\subsection{Short English statement}
Within the intended space of mildly dissipative disk diffeomorphisms, does a map lie on the boundary of the zero-topological-entropy locus exactly when it has zero entropy and is infinitely renormalizable? The space, topology, and renormalization conventions require clarification from the sources.
\subsection{Sources}
\begin{itemize}
\item Local project registry: supplied UnsolvedMath record with permanent ID \(30004181\), title ``Boundary of Zero-Entropy Dissipative Systems,'' in the Oberwolfach Reports Open Problems collection; exact source problem code \texttt{catalogue entry}. The record gives proposed year 2019 and names no individual proposer.
\item Supplied report citation: \emph{Dynamische Systeme (2020)}, Oberwolfach Reports, DOI \texttt{10.4171/owr/2019/31}, pp.\ 1855--1917. \url{https://doi.org/10.4171/owr/2019/31}. The background describes Oberwolfach Reports as documenting mathematical workshops and their open-problem discussions.
\item Supplied literature reference labeled ``Crovisier--Pujals--Tresser, zero entropy'': \url{https://arxiv.org/abs/2005.14278}.
\item Supplied literature reference labeled ``original OWR report'': \url{https://doi.org/10.4171/owr/2019/31}.
\item Literature assessment dated exactly \texttt{2026-08-22}: the supplied curation labels the problem partially solved and reports that inclusion of the zero-entropy boundary in the infinitely renormalizable class is proved for mildly dissipative disc maps, while equality and the unrestricted dissipative setting are not established. This is supplied dated curation, not independently verified. The background expressly calls it dated literature triage and requests independent verification of time-sensitive status and cited claims.
\item The registry labels the statement \texttt{corrected\_verified}; its accompanying explanation concerns preservation of mathematical content while repairing notation, typography, and sentence structure. This supplied editorial label does not establish mathematical correctness or resolution of the question.
\end{itemize}
\end{document}
